\documentclass[11pt]{article}

\usepackage[margin=1in]{geometry}
\usepackage{amsmath, amssymb, amsthm}
\usepackage{graphicx}
\usepackage{hyperref}
\usepackage{cleveref}
\usepackage{booktabs}

\graphicspath{{../figures/}}

\theoremstyle{plain}
\newtheorem{theorem}{Theorem}
\newtheorem{lemma}[theorem]{Lemma}
\newtheorem{proposition}[theorem]{Proposition}
\newtheorem{corollary}[theorem]{Corollary}
\newtheorem{conjecture}[theorem]{Conjecture}

\theoremstyle{definition}
\newtheorem{definition}[theorem]{Definition}
\newtheorem{problem}[theorem]{Problem}

\theoremstyle{remark}
\newtheorem{remark}[theorem]{Remark}

\title{Notes on Open Problems in Robot Evacuation}
\author{}
\date{\today}

\begin{document}
\maketitle

\begin{abstract}
Working notes on open problems identified in the robot evacuation literature,
with partial results and numerical evidence. This document is regenerated as
new results come in; see \texttt{CLAUDE.md} at the project root for the
current status.
\end{abstract}

\section{Open Problem Under Study}
\label{sec:problem}

We study the wireless $k$-robot evacuation problem on the unit disk at
$k=3$. Three unit-speed robots start at the centre; an exit lies at an
unknown angle $\theta$ on the boundary; once any robot's trajectory
visits the exit, a broadcast is assumed instantaneous and every robot
then walks straight to the exit. The evacuation time is the moment the
last robot reaches the exit, and the cost of an algorithm is its worst
case over $\theta$. Czyzowicz et al.~\cite{czyzowicz2014} give an
explicit algorithm~$A_3$ with UB $\frac{4\pi}{9} + \frac{2\sqrt3 + 5}{3}
+ \tfrac{1}{600} \approx 4.21930$ and matching LB $\approx 4.159$.

\section{Background}
\label{sec:background}

Algorithm $A_3$ places robots $r_1, r_2, r_3$ as follows. Fix points $A,
B$ on the boundary with $\widehat{AB}$ having length $y$ (oriented
clockwise from $A$ to $B$). Robot $r_1$ walks radially to $A$ and scans
counter-clockwise; $r_2$ walks radially to $B$ and scans clockwise; $r_3$
walks radially to $B$, scans counter-clockwise for time $y$ (so its arc
meets $r_1$'s starting point exactly), then walks straight back to the
origin (another unit of time), then walks outward to the boundary point
at angle $\pi - y/2$ (i.e.\ the midpoint of the long arc from $B$ to
$A$). The published bound takes $y_{\mathrm{paper}} = \frac{4\pi}{9} +
\frac{2\sqrt3}{3} - \frac{401}{300}$.

\section{Results So Far}
\label{sec:results}

We retain the structural skeleton of~$A_3$ and treat $y$ as a tunable
deployment parameter. Let $T(y)$ denote the worst-case evacuation time of
$A_3(y)$.

\begin{proposition}[Two special worst-case angles]
\label{prop:two-cases}
For each $y \in (0, 2\pi/3)$, define
\[
T_1(y) \;=\; 1 + \tfrac{2\pi}{3} - \tfrac{y}{2} + \sqrt{3},
\qquad
T_2(y) \;=\; 3 + y.
\]
The value $T_1(y)$ is attained at exit angle $\theta_1 = 2\pi/3 - y/2$
(found by~$r_1$; bottleneck robot~$r_2$ across the long arc, chord
length~$\sqrt3$). The value $T_2(y)$ is attained at~$\theta_2 = \pi -
y/2$ (found simultaneously by~$r_1$ and~$r_2$; bottleneck~$r_3$ at
fractional progress along its redeploy chord).
\end{proposition}

$T_1$ is strictly decreasing in~$y$ while $T_2$ is strictly increasing,
so $T_1 = T_2$ at the unique $y = \tfrac{4\pi}{9} + \tfrac{2\sqrt3}{3} -
\tfrac{4}{3}$. However, the paper's parameter is
$y_{\mathrm{paper}} = y_\star - \tfrac{1}{300}$, where $y_\star$ denotes this balance
point; the overshoot reflects a third worst-case regime.

\begin{proposition}[Third worst-case regime, numerical]
\label{prop:three-cases}
There exists a third family of worst-case angles $\theta_3(y) \in
(1+y,\; \pi - y/2)$ at which the exit is found by~$r_1$ and the
bottleneck is~$r_3$, then mid-redeploy, on a chord through the disk's
interior. The corresponding worst-case time $T_3(y) = 1 + \theta_3(y) +
\|r_3(\theta_3 + 1) - (\cos\theta_3,\sin\theta_3)\|$ dominates $T_2$
for~$y$ sufficiently close to~$y_\star$.
\end{proposition}

\begin{theorem}[Numerical improvement within the $A_3$ family]
\label{thm:main-num}
Let $y_{\mathrm{opt}} \approx 1.215858$ be the unique value in $(0,
2\pi/3)$ at which $T_1(y) = \max_\theta T_3(y, \theta)$. Then
\[
T(y_{\mathrm{opt}}) \;=\; 1 + \tfrac{2\pi}{3} - \tfrac{y_{\mathrm{opt}}}{2}
+ \sqrt{3} \;\approx\; 4.218517,
\]
improving the Czyzowicz et al.~upper bound by $4.21930 - 4.21852
\approx 0.00078 \approx 1/1280$.
\end{theorem}

At $y_{\mathrm{opt}}$ the simulator reveals four equally-binding worst-case
angles (two pairs related by the $A_3$ reflection symmetry): a
case-1 pair at $\theta \in \{1.486,\, 3.581\}$ with~$r_1$ or~$r_2$
bottleneck at chord~$\sqrt3$, and a case-3 pair at $\theta \in \{2.331,\,
2.737\}$ with~$r_3$ bottleneck at chord~$\approx 0.888$. Case~2 falls
strictly below: $T_2(y_{\mathrm{opt}}) \approx 4.2159$, slack by
$\approx 0.0026$.

\begin{theorem}[Transcendental characterization]
\label{thm:transcendental}
Let $v = \pi - y/2 - \theta_3$ and $W = \pi - 3y/2 - 1$. The critical-point
condition $\partial T_3 / \partial v = 0$ reduces to the quadratic
\[
(W-v)^2\,(2 - \sin v) - 2(W-v)\cos v + \sin v = 0,
\]
whose admissible branch is $W - v = \bigl(\cos v - \sqrt{1 - 2\sin v}\bigr)
/(2 - \sin v)$. Combining with the balance $T_1(y) = T_3(y, v)$ yields
the master equation
\begin{equation}
\label{eq:master}
\frac{\cos v + (1 - \sin v)\sqrt{1 - 2\sin v}}{2 - \sin v}
\;=\; v + \sqrt{3} - \tfrac{\pi}{3},
\end{equation}
and
\[
y_{\mathrm{opt}} = \tfrac{2}{3}\!\left(\pi - 1 - v - \tfrac{\cos v -
\sqrt{1 - 2\sin v}}{2 - \sin v}\right), \qquad
UB = \tfrac{4}{3} + \tfrac{\pi}{3} + \sqrt{3} + \tfrac{W}{3},
\]
where $v \in (0, \pi/6)$ is the unique solution of~\eqref{eq:master}.
\end{theorem}

\begin{remark}[No elementary simplification]
\label{rem:pslq}
A 40-digit PSLQ-style integer-relation search
(\texttt{FindIntegerNullVector}) over the basis $\{1,\, \pi,\, \sqrt 3,\,
\sqrt 2,\, \log 2,\, \log 3,\, \gamma\}$ finds no $\mathbb Z$-linear relation
of $v_{\mathrm{opt}}$ or the improved $UB$ with norm below~20; a subsequent
search proves no such relation of norm below~$15.2$ exists. We conclude
that $v_{\mathrm{opt}}$ (and hence $y_{\mathrm{opt}}$) is genuinely
transcendental in the usual sense: no rational-combination reduction to
elementary constants exists within the search radius.
\end{remark}

\begin{corollary}[High-precision numerical values]
\label{cor:hp}
To 40 digits,
\[
\begin{aligned}
v_{\mathrm{opt}} &= 0.20283\,77447\,40492\,87031\,05459\,53643\,95496\,52206,\\
y_{\mathrm{opt}} &= 1.21585\,78321\,29242\,90943\,42098\,32081\,15330\,17778,\\
UB              &= 4.21851\,69938\,97451\,33111\,87703\,47651\,63097\,21854.
\end{aligned}
\]
\end{corollary}

\section{Numerical Experiments}
\label{sec:experiments}

All results above are reproducible via \texttt{experiments/a3\_verify.py}
(confirms the simulator matches the paper's UB to~10 decimal places),
\texttt{experiments/a3\_optimize.py} (scalar search for $y_{\mathrm{opt}}$),
\texttt{experiments/a3\_local\_search.py} (verifies no further gain from
breaking the $L_1 = L_2$ or redeploy-angle symmetries), and
\texttt{experiments/a3\_diagnose.py} (identifies the four binding
worst-case angles at $y_{\mathrm{opt}}$).

\section{Next Steps}
\label{sec:next}

% Open directions, mirrored from CLAUDE.md.

\bibliographystyle{plain}
% \bibliography{refs}

\end{document}
